% =====================================================================
%  Преамбула отчёта о НИР (ГОСТ 7.32-2017)
%  Сборка: XeLaTeX + Biber
%  Стиль библиографии: ГОСТ Р 7.0.100-2018 (biblatex-gost)
% =====================================================================

% Кодировка и шрифты
\usepackage{fontspec}
\setmainfont{Times New Roman}
\setsansfont{Arial}
% Consolas содержит кириллицу (в отличие от Courier New) и используется
% как моноширинный шрифт; одновременно объявляется как \cyrillicfonttt
% для polyglossia, что устраняет предупреждения при наборе \verb{...}
% с русским текстом.
\setmonofont{Consolas}[Scale=0.92]
\newfontfamily\cyrillicfonttt{Consolas}[Scale=0.92]

\usepackage{polyglossia}
\setdefaultlanguage{russian}
\setotherlanguage{english}

% Геометрия страницы (ГОСТ 7.32-2017, п. 6.1):
% левое 30 мм, правое 15 мм, верхнее 20 мм, нижнее 20 мм
\usepackage[a4paper, left=30mm, right=15mm, top=20mm, bottom=20mm,
            footskip=10mm]{geometry}

% Межстрочный интервал 1.5 (ГОСТ 7.32-2017, п. 6.3)
\usepackage{setspace}
\onehalfspacing

% Абзацный отступ 1.25 см
\usepackage{indentfirst}
\setlength{\parindent}{1.25cm}

% Математика
\usepackage{amsmath, amssymb, amsthm, mathtools}

% Графика, таблицы, списки
\usepackage{graphicx}
\usepackage{booktabs}
\usepackage{longtable}
\usepackage{tabularx}
\usepackage{array}
\usepackage{multirow}
\usepackage{makecell}
\usepackage{caption}
\usepackage{subcaption}
\usepackage{float}
\usepackage{enumitem}
\setlist{nosep, leftmargin=*}

% Цвет, ссылки, листинги
\usepackage{xcolor}
\usepackage[hidelinks]{hyperref}
\hypersetup{
    pdfencoding=auto,
    bookmarksnumbered=true,
    bookmarksopen=true,
    colorlinks=false
}
\usepackage{listings}
\lstset{
    basicstyle=\ttfamily\footnotesize,
    breaklines=true,
    showstringspaces=false,
    keywordstyle=\bfseries,
    columns=fullflexible,
    captionpos=b,
    frame=single,
    aboveskip=8pt,
    belowskip=8pt
}

% Заголовки и оглавление
\usepackage{titlesec}
\usepackage{tocloft}

% Библиография по ГОСТ Р 7.0.100-2018
\usepackage[
    backend=biber,
    style=gost-numeric,
    language=auto,
    sorting=none,
    maxbibnames=4,
    minbibnames=1
]{biblatex}

% Дополнительные служебные пакеты
\usepackage{csquotes}
\usepackage{etoolbox}
\usepackage{url}
\usepackage[activate={true,nocompatibility},final,
            tracking=true,kerning=true,spacing=true,
            factor=1100,stretch=20,shrink=20]{microtype}
\usepackage{xurl}

% Разрешаем переносы внутри \verb-выражений с длинными словами
\usepackage{seqsplit}

% Разрешаем перенос длинных URL и моноширинных слов с подчёркиваниями
\PassOptionsToPackage{hyphens}{url}
\sloppy
\emergencystretch=3em
\hbadness=10000      % подавить underfull-предупреждения о хорошо
\vbadness=10000      % набранных страницах
\hfuzz=2pt           % допускать overfull <= 2pt без warning
\vfuzz=2pt

% Параметры качества разрыва строк (помогают вписать русский текст
% в узкие ячейки таблиц и абзацы со ссылками на источники)
\hyphenpenalty=1000
\tolerance=2000

% Сброс счётчиков таблиц, рисунков и формул при начале новой секции
% (требование сквозной нумерации в пределах раздела по ГОСТ 7.32-2017,
% п. 6.5.3 для рисунков, п. 6.6.2 для таблиц, п. 6.7.1 для формул)
\makeatletter
\@addtoreset{table}{section}
\@addtoreset{figure}{section}
\@addtoreset{equation}{section}
\makeatother

% =====================================================================
% Стиль ГОСТ 7.32-2017
% =====================================================================
\titleformat{\section}
    {\normalfont\fontsize{14}{17}\bfseries}
    {\thesection}{1em}{}
\titlespacing*{\section}{\parindent}{18pt}{12pt}

\titleformat{\subsection}
    {\normalfont\fontsize{14}{17}\bfseries}
    {\thesubsection}{1em}{}
\titlespacing*{\subsection}{\parindent}{14pt}{10pt}

\titleformat{\subsubsection}
    {\normalfont\fontsize{14}{17}\bfseries}
    {\thesubsubsection}{1em}{}
\titlespacing*{\subsubsection}{\parindent}{12pt}{8pt}

% Нумерация рисунков, таблиц, формул --- в пределах раздела
\renewcommand{\thefigure}{\arabic{section}.\arabic{figure}}
\renewcommand{\thetable}{\arabic{section}.\arabic{table}}
\renewcommand{\theequation}{\arabic{section}.\arabic{equation}}

% Подписи по ГОСТ
\captionsetup[figure]{
    labelsep=endash, name=Рисунок,
    justification=centering, singlelinecheck=true,
    font=normalsize
}
\captionsetup[table]{
    labelsep=endash, name=Таблица,
    justification=raggedright, singlelinecheck=false,
    font=normalsize
}
% Подписи внутри longtable наследуют те же настройки --- название
% «Таблица N — ...» помещается над таблицей по левому краю (ГОСТ 7.32-2017).
\captionsetup[longtable]{
    labelsep=endash, name=Таблица,
    justification=raggedright, singlelinecheck=false,
    font=normalsize
}

% Колонтитулы
\usepackage{fancyhdr}
\pagestyle{fancy}
\fancyhf{}
\fancyfoot[C]{\thepage}
\renewcommand{\headrulewidth}{0pt}
\renewcommand{\footrulewidth}{0pt}

\renewcommand{\contentsname}{\hfill СОДЕРЖАНИЕ\hfill}

% =====================================================================
% Пользовательские команды
% =====================================================================
\newcommand{\figref}[1]{рисунок~\ref{#1}}
\newcommand{\tabref}[1]{таблица~\ref{#1}}
\newcommand{\secref}[1]{раздел~\ref{#1}}
\newcommand{\eng}[1]{\textit{#1}}

\DeclareMathOperator*{\argmax}{arg\,max}
\DeclareMathOperator*{\argmin}{arg\,min}
\DeclareMathOperator{\softmax}{softmax}
\DeclareMathOperator{\KL}{KL}
\DeclareMathOperator{\MMD}{MMD}
\DeclareMathOperator{\PSI}{PSI}

% Метка последней страницы для реферата
\usepackage{lastpage}
